\subsection{Настройка движка Camunda под нагрузку}

\subsubsection{Задача}
\begin{itemize}
  \item Настроен встроенный брокер Camunda для эффективной работы при высокой нагрузке с external tasks.
  \item Job Executors оптимизированы под асинхронное выполнение задач с минимальным временем ожидания.
  \item Запись истории в слой персистенции отключена; при этом логика работы брокера полностью сохраняется.
\end{itemize}

\subsubsection{Аналитика задачи}

Потребоность в задаче появилась после нагрузочных тестов, которые выявили некоторые промедления при взаимодейтсвии с системой бизнес процессов.
25.03.2026 после окончательного аудита и профилирования сервисов, задача получила высший приоритет, поэтому я решил её взять, чтобы научиться конфигурировать
готовые системы под тяжёлую нагрузку. Для постановки задачи мы собрались с системным аналитиком и архитектором.
Обнаружилось, что камунда при реализации шаблона $external tasks$, использует свой встроенный брокер сообщений для делегирования исполнения кода на сторону бизнес логики.
Я заметил, что есть некоторые задержки при исполнении сервисных задач, при этом задержка происходит до исполнения задачи на стороне бизнес логики, что являлось показателем больших таймаутов на опрос топиков с задачами.
Для более быстрой обработки нужно найти и отрегулировать конфигурацию движка Camunda. Далее само выполнение и забор сервисных задач выполнялся без многопоточной обработки, так следует из документации [5].
Аналитик и архитектор также подметили, что для разгрузки слоя персистенции можно отключить запись истории в таблицы $postgres$, так как это не входит в перечень обязательного аудита, а сами исторические таблицы не нужны для работы процессов, так как есть операционные таблицы, в которых хранятся динамические данные процесса, поэтому было решено изначально выключить историю, а в будущем, перенаправить её в отдельную базу данных.

\subsubsection{Написание конфигурации Job Executor}

Как я вычитал из документации -- по умолчанию движок запускает $Job Executor$ с минимальными параметрами, ориентированными на однопоточное выполнение.

Для настройки я открыл файл \texttt{bpn-platform.xml}, который копируется в контейнер и применяется при запуске движка Camunda. В этом файле я добавил следующие параметры:
параметр $core-pool-size$ задаёт базовое число потоков в пуле -- я установил значение 3, чтобы движок постоянно держал несколько потоков готовыми к захвату задач;
параметр $max-pool-size$ ограничивает максимальный рост пула при пиковой нагрузке -- выставлено значение 10, которое пропорционально потокам процессора, которые выделены в рамках лимита $K8S$ для пода Camunda;
параметр $max-jobs-per-acquisition$ определяет, сколько задач движок захватывает за один цикл опроса; я сделал увеличение с 1 до 10, что позволит обрабатывать пачки задач за один обход без многократных обращений к брокеру;
Для снижения задержек так же оказался полезен параметр $wait-time-in-millis$, так как уменьшение интервала ожидания между циклами опроса до 300~мс позволило существенно сократить время от появления задачи в очереди до её захвата воркером.
Параметр $lock-time-in-millis$ оставлен достаточно большим (300000~мс), чтобы задача не освобождалась повторно во время длительной обработки.

Полный листинг конфигурации вынесен в приложение~\ref{appendix:camunda-config}.

\subsubsection{Настройка параметров длительного опроса}

Клиенты внешних задач периодически опрашивают REST API движка на наличие новых заданий по своим топикам.
При коротком тайм-ауте клиент получает пустой ответ и в нашей системе, так как реализована система повторов, то вынужден тут же отправлять следующий запрос, создавая избыточную нагрузку на брокер.

Для устранения этого я настроил механизм длительного опроса: установил параметр $async-response-timeout$ равным 10000мс.
При данной настройке движок удерживает соединение до появления задачи или до истечения тайм-аута, после чего возвращает ответ. После того, как я удостоверился при помощи инструментального профилирования с замерами времени выполнения методов и сбором сквозной истории, что клиенты начали получать задачи быстрее.
Далее я выставил параметр lock-duration в 20000мс.

\subsubsection{Отключение уровня хранения истории}

Camunda поддерживает несколько уровней хранения истории: full, audit, activity, none.
Как я убедился, по умолчанию уровень установлен в full, при котором движок записывает в исторические таблицы персистенцию --- переходы по шагам, значения переменных, временные метки.
Эти записи выполняются синхронно в рамках той же транзакции, что и продвижение токена. Как раз в этом случае мы и можем стать узким местом, или ещё что хуже -- перегружать базу историей.

Поскольку исторические таблицы не используются ни для мониторинга, ни для аудита в рамках текущего контура, я установил параметр none.
При данном значении движок полностью пропускает запись в исторические таблицы, не затрагивая операционные таблицы (\texttt{ACT\_RU\_*}).

\subsubsection{Конфигурационные файлы}

Итоговые конфигурационные файлы, применённые в рамках данной задачи, вынесены в приложение~\ref{appendix:camunda-config}.
Данный конфиг подключается после сборки контейнера с плагинами, он передаётся в контейнер с образом camunda, чтобы при старте она самогла применить не стандартные настройки.

\subsubsection{Проверка быстрой асинхронной работы job executors}

Для проверки корректности настройки Job Executor я запустил несколько экземпляров процессов, содержащих сервисные задачи, которые как раз реализуют external tasks.
Во время профилирования было заметно, как задачи стали захватываться именно пакетами и распределяться по нескольким потокам из пула.

\subsubsection{Проверка отсутствия истории в слое персистенции при движении по схеме}

После отключения истории я запустил несколько экземпляров процессов (из проверки предыдущих шагов) и выполнил прямые запросы к базе данных PostgreSQL на исторические таблицы \texttt{ACT\_HI\_PROCINST}, \texttt{ACT\_HI\_ACTINST} и \texttt{ACT\_HI\_TASKINST}.
Во всех таблицах отсутствовали новые записи, соответствующие запущенным экземплярам, что подтвердило отключение записи истории.
Одновременно я проверил также динамические таблицы: там , как я увидел были всё те же операционные данные активных экземпляров. Процессы создавались без проблем и жили без аварийных завершений

\subsubsection{Вывод}

В рамках задачи я настроил движок Camunda для работы с высоконагруженной системой.
Настроил Job Executor за счёт увеличения пула потоков и сокращения интервала опроса устранила задержки при захвате сервисных задач.
Включил механизм длительного опроса для клиентов external tasks, что снизило число избыточных запросов к брокеру.
Отключил уровень хранения истории и тем самым разгрузил слой персистенции без нарушения логики исполнения процессов.
